\documentclass[a4paper,12pt]{article}
\usepackage{amsmath}
\usepackage{amssymb}
\usepackage{booktabs}
\usepackage{geometry}
\usepackage{graphicx}
\usepackage[portuguese]{babel}  
% Pacote para códigos em Python
\usepackage{listings}
\usepackage{xcolor}
\renewcommand{\lstlistingname}{Algoritmo}
% Configuração do ambiente listings para Python
\lstset{
    language=Python,
    basicstyle=\ttfamily\small,
    keywordstyle=\color{blue}\bfseries,
    stringstyle=\color{red},
    commentstyle=\color{gray}\itshape,
    numbers=left,
    numberstyle=\tiny\color{gray},
    stepnumber=1,
    numbersep=5pt,
    backgroundcolor=\color{white},
    frame=single,
    tabsize=4,
    showspaces=false,
    showstringspaces=false,
    breaklines=true,
    breakatwhitespace=true,
    captionpos=b
}
\geometry{margin=0.5in}

\title{Lista Mínimos Quadrados - Prática}
\author{}
\date{}

\begin{document}
\maketitle

\begin{enumerate}
    \item Para cada conjunto de dados nos arquivos \texttt{'ex1\_<item>.npz'}, com os itens variando de (a) até (d), leia os arquivos e aproxime os dados através de um polinômio. Justifique a escolha do grau do polinômio usando o plot dos dados e do modelo encontrado.
    
    \begin{lstlisting}[gobble=8,caption=Exemplo de leitura de arquivo:]
        dados = np.load('ex1_a.npz')
        x = dados['x']
        y = dados['y']
    \end{lstlisting}

    \item \textbf{Modelo polinomial com múltiplas variáveis.}  
    A ideia de modelos polinomiais pode ser estendida a partir do caso em que há apenas uma variável (feature).  
    Neste exercício, consideramos um modelo quadrático (de grau dois) com 2 variáveis, isto é, $x$ é um vetor em $\mathbb{R}^2$.  
    O modelo tem a forma
    \[
    \hat{f}(x) = a + b_1 x_1 + b_2 x_2 
    + c_1 x_1^2 + c_2 x_2^2
    + c_3 x_1 x_2
    \]
    onde o escalar $a$, o vetor $b \in \mathbb{R}^2$ e o vetor $c \in \mathbb{R}^3$ são, respectivamente, os coeficientes de ordem zero, primeira e segunda do modelo.

    \begin{enumerate}
        \item Leia os dados do arquivo  \texttt{'ex2.npz'}, que contem valores de $x$, $y$ e $z$ e encontre uma aproximação para $f(x,y)=z$.
        \item Calcule o $\mathbf{\text{rms}}$ entre o modelo encontrado e os dados.
        \item Plote a curva encontrada pelo seu modelo.
    \begin{lstlisting}[gobble=8,caption=Exemplo de plot de superficies]
        fig = plt.figure()
        ax = fig.add_subplot(projection='3d')
        surf = ax.plot_surface(X, Y, z, cmap='viridis', edgecolor='none', antialiased=True)
        fig.colorbar(surf, ax=ax, shrink=0.6)
        ax.set_xlabel('x')
        ax.set_ylabel('y')
        ax.set_zlabel('z')
        ax.set_title('3D surface plot of f(x,y)')
    \end{lstlisting}


    \end{enumerate} 


    \item O conjunto de dados contém o número de transistores $N$ em 13 microprocessadores, e seu ano de lançamento.
    
    \begin{table}[htb!]
        \centering
        \begin{tabular}{cr}
        \hline
        Ano  & \multicolumn{1}{c}{Transistores} \\ \hline
        1971 & 2,250                            \\
        1972 & 2,500                            \\
        1974 & 5,000                            \\
        1978 & 29,000                           \\
        1982 & 120,000                          \\
        1985 & 275,000                          \\
        1989 & 1,180,000                        \\
        1993 & 3,100,000                        \\
        1997 & 7,500,000                        \\
        1999 & 24,000,000                       \\
        2000 & 42,000,000                       \\
        2002 & 220,000,000                      \\
        2003 & 410,000,000                      \\ \hline
        \end{tabular}
    \end{table}
    \begin{enumerate}
        \item Encontre a melhor aproximação de mínimos quadrados para aproximar os dados. Justifique a escolha do modelo escolhido para aproximar os dados.
        \item Calcule o $\mathbf{\text{rms}}$ do erro da aproximação nos dados. Ploete o modelo que você encontrou junto com o conjunto de dados. \textit{Dica}: Utilize \texttt{plt.yscale('log')} para melhor visualização.
        \item Use o seu modelo para prever o número de transistores em um microprocessador introduzido em 2015. Compare essa predição com o microprocessador IBM Z13, lançado em 2015, que possui cerca de $4 \times 10^9$ transistores.

        \item Compare o seu resultado com a lei de Moore, que afirma que o número de transistores por circuito integrado aproximadamente dobra a cada um ano e meio a dois anos.
\end{enumerate}
    \item Um laboratório está estudando a degradação da capacidade de baterias de íons de lítio utilizada em veículos elétricos. 
Para diferentes condições de uso, foram registradas medições ao longo do número de ciclos de carga/descarga.

Os dados estão armazenados no arquivo \texttt{ex4.npz}, fornecido junto com esta lista. Os dados do arquivo são:

\begin{center}
\begin{tabular}{l l}
\texttt{cycle}        & número de ciclos de carga/descarga (inteiro positivo) \\
\texttt{temperature}  & temperatura média de operação durante o ciclo (em $^\circ$C) \\
\texttt{capacity}     & capacidade medida da bateria (em \% da capacidade nominal) \\
\end{tabular}
\end{center}

Queremos ajustar um modelo que relacione a capacidade da bateria com o número de ciclos e a temperatura, da forma
\[
\hat{C}(n, T) 
= \theta_0 
+ \theta_1 n 
+ \theta_2 T 
+ \theta_3 n^2 
+ \theta_4 T^2 
+ \theta_5 nT,
\]
onde $n$ é o número de ciclos, $T$ é a temperatura (em $^\circ$C), e os coeficientes $\theta_0,\dots,\theta_5$ serão obtidos por mínimos quadrados.

\begin{enumerate}

    \item Escreva o problema de ajuste em mínimos quadrados na forma matricial
    \[
        A\theta \approx C,
    \]
    especificando explicitamente a matriz de projeto $A \in \mathbb{R}^{N \times 6}$, o vetor de parâmetros $\theta \in \mathbb{R}^6$ e o vetor de observações $C \in \mathbb{R}^N$.


    \item Calcule numericamente $\hat{\theta}$. Em seguida, calcule o erro RMS do ajuste:

    \item Use o modelo ajustado para prever a capacidade da bateria nas seguintes situações:
    \begin{enumerate}
        \item $n = 500$ ciclos, $T = 25^\circ$C;
        \item $n = 500$ ciclos, $T = 45^\circ$C;
        \item $n = 1500$ ciclos, $T = 25^\circ$C.
    \end{enumerate}
    Comente o efeito da temperatura e do número de ciclos sobre a capacidade prevista.

    \item Plote, em um mesmo gráfico, os valores medidos de capacidade versus número de ciclos para uma faixa estreita de temperatura (por exemplo, $T$ entre $23^\circ$C e $27^\circ$C) e a curva correspondente dada pelo modelo $\hat{C}(n,T)$ fixando $T = 25^\circ$C. 
    Com base nesse gráfico, discuta qualitativamente se o modelo quadrático em $n$ parece adequado.

\end{enumerate}

\end{enumerate}
\end{document}